\section{Rowing-Specific and Experience Measures}
\label{sec:study_design_measures_rowing_specific}

The established embodiment items will be complemented by custom rowing-specific items. These items are necessary because the main thesis comparison is not only whether the virtual body feels owned or controlled, but also whether the representation feels like rowing and whether the movement mapping feels believable.

The custom per-condition items will cover:

\begin{itemize}
    \item perceived match between real movement and virtual movement,
    \item perceived timing or responsiveness,
    \item rowing realism or sport-action plausibility,
    \item perceived connection between physical effort and virtual rowing outcome,
    \item comfort, discomfort, or noticeable mismatch.
\end{itemize}

These custom items will be treated as exploratory thesis-specific measures. They will not be mixed into VEQ Ownership or VEQ Agency scores.

After both conditions, participants will answer comparative preference questions. These will ask which representation they preferred overall and which representation felt more convincing for rowing. Open comments will be collected to identify mismatch sources, comfort issues, and details that are not captured by numeric ratings.

\drafttodo{TODO:CONFIRM final custom item list and response scale before pilot testing.}

